% Source printed pp. 420-422 / PDF pp. 428-430.
\section{Cartan 分解与不稳定化定理}\label{src:8.3}
\begin{epigraph}
我已经确信, 神学也有它的优点.
\epigraphauthor{Paul Gordan}
\end{epigraph}
Cartan 分解 (\bookref{8.3.1}) 是证明不稳定化定理 (\bookref{8.3.6}) 的关键代数工具; 后者将在下一节用于证明 Hilbert--Mumford 判据 (\bookref{8.4.4}).
\subsection{Cartan 分解}\label{src:8.3.1-subsection}
\textit{Cartan 分解}---有时称为 \textit{Iwahori 分解}或 \textit{Cartan--Iwahori--Matsumoto 分解}---断言: 对 DVR $R$ 的分式域 $K$ 上的元素 $g\in G(K)$, 可在左右两侧乘以 $G(R)$ 中的元素, 使结果来自一个单参数子群. 对单参数子群 $\lambda:\mathbb G_m\to G$, 用 $\lambda|_K\in G(K)$ 表示复合
\[
\begin{tikzcd}[ampersand replacement=\&]
\Spec K\arrow[r]\&\mathbb G_m\arrow[r,"\lambda"]\&G
\end{tikzcd}
\]
的像, 其中第一个映射由将 $t$ 映到 $R$ 中一致化元的 $\Bbbk$-代数映射 $\Bbbk[t]_t\to K$ 定义.
\begin{theorem}[Cartan 分解]\label{src:8.3.1}
设 $G$ 是代数闭域 $\Bbbk$ 上的约化群, $R$ 是 $\Bbbk$ 上剩余域为 $\Bbbk$, 分式域为 $K$ 的完备 DVR. 则对每个 $g\in G(K)$, 存在 $h_1,h_2\in G(R)$ 和单参数子群 $\lambda:\mathbb G_m\to G$, 使
\[
g=h_1\lambda|_Kh_2.
\]
\end{theorem}
\begin{proof}
我们已在命题~\bookref{7.10.30} 中证明此结论, 并证明 Cartan 分解的存在性刻画了约化性. 经典证明见 \cite[推论 2.17]{IM65}, \cite[定理 2.1]{Ses72} 和 \cite[\S4]{BT72}.
\end{proof}
\begin{remark}[等价表述]\label{src:8.3.2}
设 $T\subseteq G$ 是极大环面. 上述定理等价于等式
\[
G(K)=G(R)T(K)G(R).
\]
为看出分解如何推出此等式, 选择 $h\in G(R)$ 使 $h\lambda|_Kh^{-1}\in T(K)$. 则
\[
g=h_1\lambda|_Kh_2=\underbrace{(h_1h^{-1})}_{\in G(R)}\underbrace{(h\lambda|_Kh^{-1})}_{\in T(K)}\underbrace{(hh_2)}_{\in G(R)}.
\]
反之, 设 $g=h_1th_2$, 其中 $h_1,h_2\in G(R)$, $t\in T(K)$. 若写 $T\cong\mathbb G_m^r$, 并以 $\pi\in R$ 为一致化参数, 则 $t=(u_1\pi^{d_1},\ldots,u_r\pi^{d_r})$, 其中 $u_i\in R^\times$, $d_i\in\mathbb Z$. 将 $h_1$ 替换为 $h_1\cdot(u_1,\ldots,u_r)$ 后, 可写 $g=h_1\lambda|_Kh_2$, 其中 $\lambda:\mathbb G_m\to T\subseteq G$ 是由 $t\mapsto(t^{d_1},\ldots,t^{d_r})$ 给出的单参数子群.
\end{remark}
\begin{remark}[\texorpdfstring{$\mathrm{GL}_n$}{GLn} 的情形]\label{src:8.3.3}
$\mathrm{GL}_n$ 的 Cartan 分解可用初等线性代数证明. 设 $g=(g_{ij})\in\mathrm{GL}_n(K)$. 作行列变换后, 可设 $g_{1,1}=\pi^d$ 在所有 $g_{ij}$ 中赋值最小, 其中 $\pi\in R$ 是一致化元. 对每个 $k\geq2$, 写 $g_{k,1}=u\pi^e$. 现在作行变换, 将第 $n$ 行 $r_n$ 替换为 $r_n-u\pi^{e-d}r_1$. 如此可使对 $k\geq2$ 有 $g_{k,1}=0$. 作类似的列变换还可使 $g_{1,k}=0$ ($k\geq2$). 结论遂由归纳法得到.
\end{remark}
\begin{exercise}\label{src:8.3.4}
\begin{enumerate}[(a)]
\item 设 $X\subseteq\mathbb P(V)$ 是 $\mathbb G_m$-等变局部闭子概形, 其中 $V$ 是有限维 $\mathbb G_m$-表示. 证明 $[X/\mathbb G_m]$ 分离当且仅当 $X$ 没有 $\mathbb G_m$-不动点. 换言之, 对角 $[X/\mathbb G_m]\to[X/\mathbb G_m]\times[X/\mathbb G_m]$ 有限当且仅当它拟有限.
\item 设约化群 $G$ 作用于 $\Bbbk$ 上代数空间 $X$. 证明 $[X/G]$ 分离当且仅当对每个单参数子群 $\lambda:\mathbb G_m\to G$, 诱导商叠 $[X/\mathbb G_m]$ 分离.
\end{enumerate}
\textit{提示:} 应用 Cartan 分解验证赋值判别法.
\end{exercise}
\begin{remark}\label{src:8.3.5}
若将 $\mathbb G_m$ 替换为一般线性约化群, 则习题~\bookref{8.3.4}\bookref{8.3.4} 不成立; $\mathrm{SL}_2$ 在拟仿射簇上的自由作用给出反例 (见习题~\bookref{4.9.4}(d)).
\end{remark}
\subsection{不稳定化定理}\label{src:8.3.2-subsection}
\begin{theorem}[不稳定化定理]\label{src:8.3.6}
设代数闭域 $\Bbbk$ 上的约化群 $G$ 作用于 $\Bbbk$ 上有限型仿射概形 $X$. 给定 $x\in X(\Bbbk)$, 存在单参数子群 $\lambda:\mathbb G_m\to G$, 使 $x_0:=\lim_{t\to0}\lambda(t)\cdot x$ 存在且其 $G$-轨道闭.
\end{theorem}
\begin{proof}
令 $R=\Bbbk[\![t]\!]$, 分式域为 $K=\Bbbk((t))$. 可选取元素 $g\in G(K)$ 和交换图
\[
\begin{tikzcd}[ampersand replacement=\&]
\Spec K\arrow[r]\arrow[d,hook]\&Gx\arrow[d,hook]\\
\Spec R\arrow[r,"\widetilde g"']\&X,
\end{tikzcd}
\]
其中上方映射由复合 $\Spec K\xrightarrow{g}G\to Gx$ 给出, 且 $y:=\widetilde g(0)$ 的 $G$-轨道闭. 由 Cartan 分解 (\bookref{8.3.1}), 存在 $h_1,h_2\in G(R)$ 和单参数子群 $\lambda:\mathbb G_m\to G$, 使 $h_1g=\lambda|_Kh_2$. 对 $a\in G(R)$ 和 $b\in X(R)$, 一般有 $(a\cdot b)(0)=a(0)\cdot b(0)$. 将此应用于 $h_1\in G(R)$ 和 $\widetilde g\in X(R)$, 得到
\begin{align*}
\lim_{t\to0}\lambda(t)h_2(t)\cdot x
&=\lim_{t\to0}h_1(t)g(t)\cdot x\\
&=h_1(0)\cdot\widetilde g(0)=h_1(0)\cdot y\in Gy.
\end{align*}
我们断言, 相关但可能不同的极限 $\lim_{t\to0}\lambda(t)h_2(0)\cdot x$ 存在, 且也属于闭轨道 $Gy$. 证明此断言后, 使用单参数子群 $h_2(0)^{-1}\lambda h_2(0)$ 即得到定理:
\[
\lim_{t\to0}(h_2(0)^{-1}\lambda h_2(0))(t)\cdot x
=h_2^{-1}(0)\cdot\lim_{t\to0}\lambda(t)h_2(0)\cdot x\in Gy.
\]
首先, 为证明 $\lim_{t\to0}\lambda(t)h_2(0)\cdot x$ 存在, 可用下面的命题~\bookref{B.1.21}(1) 归约到 $X=\mathbb A(V)$ 是 $G$-表示的情形. 可选取 $V\cong\Bbbk^n$ 的一组基, 使 $\lambda$-作用的权为 $\lambda_1,\ldots,\lambda_n$. 又可写 $h_2\cdot x=(a_1,\ldots,a_n)\in X(R)$, 其中每个 $a_i\in\Bbbk[\![t]\!]$, 并分解 $a_i=a_i(0)+a_i'$, $a_i'\in(t)$. 因为
\begin{equation}\label{src:8.3.7}
\begin{aligned}
\lim_{t\to0}\lambda(t)h_2(t)\cdot x
&=\lim_{t\to0}\bigl(t^{\lambda_1}(a_1(0)+a_1'),\ldots,\\
&\hspace{4.5em}t^{\lambda_n}(a_n(0)+a_n')\bigr)
\end{aligned}
\end{equation}
存在, 对每个 $\lambda_i<0$ 的 $i$ 必有 $a_i(0)=0$, 进而推出 $\lim_{t\to0}\lambda(t)h_2(0)\cdot x$ 存在.

最后, 为证明此极限属于 $Gy$, 可用命题~\bookref{B.1.21}(2) 得到 $G$-等变映射 $f:X\to\mathbb A(W)$, 满足 $f^{-1}(0)=Gy$. 于是只须证明 $\lim_{t\to0}\lambda(t)h_2(0)\cdot f(x)=0$. 计算 \bookeqref{8.3.7} 中的极限, 与上面相同的论证表明, $\lim_{t\to0}\lambda(t)h_2(t)\cdot f(x)=0$ 推出 $\lim_{t\to0}\lambda(t)h_2(0)\cdot f(x)=0$. 另见 \cite[第 53 页]{GIT} 和 \cite[定理 1.4]{Kem78}.
\end{proof}
\begin{corollary}[不稳定化定理 II]\label{src:8.3.8}
设 $\mathcal X$ 是代数闭域 $\Bbbk$ 上具有仿射对角的有限型代数叠. 设 $x\leadsto x_0$ 是 $\Bbbk$-点的特殊化, 且稳定子 $G_{x_0}$ 线性约化. 则存在表示特殊化 $x\leadsto x_0$ 的态射 $[\mathbb A^1/\mathbb G_m]\to\mathcal X$.
\end{corollary}
\begin{proof}
局部结构定理 (\bookref{7.7.1}) 给出平展态射 $[\Spec A/G_{x_0}]\to\mathcal X$ 以及映到 $x_0$ 的点 $w_0$. 必要时将 $\Spec A$ 替换为 $G_{x_0}$-不变仿射子概形后, 可设 $w_0$ 是闭点. 特殊化 $x\leadsto x_0$ 提升为 $[\Spec A/G_{x_0}]$ 中的特殊化 $w\leadsto w_0$, 可选取 $w$ 所对应轨道的代表 $\widetilde w\in\Spec A$. 不稳定化定理 (\bookref{8.3.6}) 给出单参数子群 $\lambda:\mathbb G_m\to G$, 使 $\widetilde w_0=\lim_{t\to0}\lambda(t)\cdot\widetilde w$ 存在且其轨道闭. 由仿射 GIT (\bookref{7.5.13}), $\overline{G\widetilde w}$ 中有唯一闭轨道, 所以 $\widetilde w_0\in\Spec A$ 映到 $w_0$. $t\mapsto\lambda(t)\cdot\widetilde w$ 的 $\mathbb G_m$-等变延拓 $\mathbb A^1\to X$ 定义代数叠态射 $[\mathbb A^1/\mathbb G_m]\to[\Spec A/G_{x_0}]$, 使特殊化 $1\leadsto0$ 的像是 $w\leadsto w_0$. 复合 $[\mathbb A^1/\mathbb G_m]\to[\Spec A/G_{x_0}]\to\mathcal X$ 即为所求映射.
\end{proof}
